Initial results for the calculation were found using a semi-micro cuvette and after many trials and measurements all produced physically impossible values, particularly, refractive indices less than 1 which is the value for vacuum that cannot be exceeded. Many adjustments were made to the alignment of the setup and the measurements and calculations were redone and tripple-checked without much improvement, finally the plastic semi-micro cuvette was swapped for a quartz macro cuvette. Instant improvements were observed in the results and it was determined that the protruding edges of the semi-micro cuvette were clipping the first order fringe beams exiting the liquid. This resulted in errors in all of the calculations as the beams were diverted from their natural diffraction angles.

Given that this experiment was able to produce results for miniscule changes in refractive index as is evident by the ability of the setup to differentiate between water and ethanol which have very similar refractive indices, it is possible the theory could be expanded upon to determine impurities. Based on the observations and theoretical calculations, the LDR should be able to detect a change of at least $5 \times 10^{-4}$ in the refractive index, this would correlate to an ability to detect a roughly 2\% impurity in a substance. While this wouldn't qualify as lab or weapons grade purity, simple modifications could be made to increase the sensitivity to detect a change in refractive index up to $1 \times 10^{-6}$ and this process could be used to determine very minute impurities in any transparent or translucent liquids and likely gases.

